\documentclass[10pt,a4paper]{article}
\usepackage[T1]{fontenc}
\usepackage[margin=2cm]{geometry}
\usepackage{isabelle,isabellesym}
\usepackage{amsmath,amssymb}
\usepackage{pdfsetup}

\urlstyle{rm}
\isabellestyle{it}

\begin{document}

\title{Admissibility Lemmas for the Limited Augmented Zarankiewicz Number}
\author{Carles Mar\'in}
\maketitle

\begin{abstract}
The limited augmented Zarankiewicz number of Qi, Cui and Xu~\cite{QCX-BSR,QCX-L} counts the cells of a maximum
$C_4$-free set in an $m\times n$ grid together with pairs of further cells that satisfy two local conditions, (C2)
and (C3). We formalize the combinatorial lemmas that make constructions and upper bounds work: a \emph{label lemma}
(if every occupied support lies in a label, any three occupied columns of a row form a triple of a given shadow, and
the pair is compatible with the labels, then the pair is admissible), the \emph{blocks lemma} (pairs arranged
around a triangle of three column blocks whose face stays empty are admissible), the lemma on vertical pairs and the
lemma that at most one row is full for maximum bases with $m\ge\binom n2$, for every $n$; the reformulation of (C2)
and (C3) as distance at least three in the graph of occupied rectangles; the bound $2|E_2|\le(n-2)m-\binom n2+1$ on
the number of pairs; and the arithmetic of the linear gap between the weak and the original numbers in three columns.
\end{abstract}

\paragraph{Acknowledgement.} The formalization and the text of this document were prepared with the
assistance of generative AI systems, and checked by the author.

\tableofcontents

\section*{Introduction}

A \emph{base} is a $C_4$-free set $E_1$ of cells of maximum size $z(m,n)$; a family $E_2$ of disjoint pairs of
further cells is admissible when every pair satisfies (C2), its two opposite cells are not both occupied, and (C3),
no further row and column close the five-cell pattern with the pair. The theory \texttt{ZL\_Labels} proves the label
lemma, the blocks lemma, which is Lebedev's three-block mechanism~\cite{Leb} in the form needed for a fixed number of
columns, and, in a locale describing maximum bases in the regime $m\ge\binom n2$, the lemmas on vertical pairs and on
full rows. The theory \texttt{ZL\_Counts} proves that (C2) and (C3) together say that the two cells of a pair are at
distance at least three in the graph whose edges are the opposite corners of occupied rectangles, the endpoint bound,
and the arithmetic of the weak gap.

\parindent 0pt\parskip 0.5ex

\input{session}

\bibliographystyle{abbrv}
\bibliography{root}

\end{document}
